\section*{Bloque D. Ampliación: Teoría y Aplicaciones Prácticas (70\% / 30\%)}

\noindent\textit{Preguntas 91 a 180. Ampliación del banco oficial con 63 preguntas teóricas y conceptuales (91 a 153) y 27 aplicaciones cuantitativas prácticas (154 a 180).}

\subsection*{Parte I. Preguntas teóricas y conceptuales (91 a 153)}

\q{En el enfoque de conjuntos de producción $Y \subset \mathbb{R}^m$, la propiedad de libre disposición (monotonicidad) establece que:}
\begin{opts}
    \item es posible producir cantidades positivas sin emplear ningún insumo físico.
    \item toda combinación convexa de dos planes factibles requiere igual capital.
    \item si un plan $y \in Y$, cualquier plan menos eficiente también es factible.
    \item el producto marginal de cada insumo es estrictamente creciente en $Y$.
\end{opts}

\q{La propiedad de inacción ($0 \in Y$) junto con la convexidad del conjunto de producción $Y$ implica económicamente que la tecnología exhibe:}
\begin{opts}
    \item rendimientos a escala no crecientes (constantes o decrecientes).
    \item rendimientos a escala estrictamente crecientes en todo el dominio.
    \item costos fijos irrecuperables en el horizonte de largo plazo.
    \item elasticidad de sustitución infinita entre todos los factores.
\end{opts}

\q{En la producción a corto plazo con capital fijo $\bar{K}$ y trabajo variable $L$, la ley de rendimientos marginales decrecientes establece que a partir de cierto nivel de $L$:}
\begin{opts}
    \item el producto total $PT$ necesariamente se vuelve negativo de inmediato.
    \item los rendimientos a escala de largo plazo pasan a ser decrecientes.
    \item el costo marginal de corto plazo disminuye conforme crece el producto.
    \item el producto marginal del trabajo $PMg_L$ disminuye al añadir más $L$.
\end{opts}

\q{En la teoría clásica de las tres etapas de producción a corto plazo, la etapa II (región económica relevante para una firma racional) se define como el intervalo donde:}
\begin{opts}
    \item $PMg_L$ es creciente y supera estrictamente al $PMe_L$.
    \item $PMe_L$ es decreciente y $0 \le PMg_L \le PMe_L$.
    \item $PMg_L$ es negativo mientras el $PMe_L$ alcanza su punto máximo.
    \item $PMe_L$ es creciente y la elasticidad producto-trabajo supera a $1$.
\end{opts}

\q{Cuando una firma opera exactamente en el punto donde el producto medio del trabajo ($PMe_L$) alcanza su máximo, la elasticidad del producto respecto al trabajo ($\varepsilon_{Q,L}$) es igual a:}
\begin{opts}
    \item cero ($0$), pues el producto total alcanza su máximo.
    \item un valor estrictamente mayor que uno en toda la etapa II.
    \item el cociente negativo entre $PMg_L$ y el producto medio.
    \item la unidad ($1$), pues $PMg_L = PMe_L$.
\end{opts}

\q{En una función de producción con rendimientos constantes a escala, si la firma opera en la etapa I respecto al trabajo ($PMg_L > PMe_L$), entonces el producto marginal del capital ($PMg_K$) es:}
\begin{opts}
    \item estrictamente positivo e igual al producto medio del capital.
    \item creciente y superior al costo de uso del capital $r$.
    \item estrictamente negativo, revelando ineficiencia técnica.
    \item idénticamente igual a la tasa marginal de sustitución técnica.
\end{opts}

\q{La relación entre el producto marginal ($PMg_L$) y el producto medio ($PMe_L$) implica que si $PMg_L < PMe_L$ en un rango de producción, entonces necesariamente:}
\begin{opts}
    \item $PMe_L$ está disminuyendo al aumentar la cantidad de trabajo $L$.
    \item $PMe_L$ está aumentando a una tasa decreciente con el trabajo $L$.
    \item $PMg_L$ debe ser necesariamente negativo en todo ese intervalo.
    \item $PT$ ha alcanzado ya su punto de inflexión y comienza a caer.
\end{opts}

\q{La convexidad estricta de las isocuantas de una función de producción diferenciable $Q = f(K,L)$ es matemáticamente equivalente a que:}
\begin{opts}
    \item la función de producción sea homogénea de grado uno en $(K,L)$.
    \item la $TMST_{L,K}$ sea decreciente a lo largo de la isocuanta.
    \item los productos marginales de ambos factores sean constantes.
    \item la elasticidad de sustitución $\sigma$ sea estrictamente igual a cero.
\end{opts}

\q{La elasticidad de sustitución técnica $\sigma$ entre capital y trabajo mide formalmente:}
\begin{opts}
    \item el cambio porcentual en $Q$ por cada $1\%$ de aumento en todos los insumos.
    \item la pendiente absoluta de la senda de expansión del producto de largo plazo.
    \item el cambio porcentual en $K/L$ por cada $1\%$ de cambio en la $TMST_{L,K}$.
    \item el cociente entre la productividad media del capital y la del trabajo.
\end{opts}

\q{Para la función de producción Cobb-Douglas $Q = A K^\alpha L^\beta$ con $\alpha, \beta > 0$, la elasticidad de sustitución $\sigma$ es siempre:}
\begin{opts}
    \item igual a $\alpha + \beta$, dependiendo de los rendimientos a escala.
    \item igual a $\alpha/\beta$, determinada por la intensidad relativa de capital.
    \item igual a $0$, porque los factores guardan proporciones técnicas fijas.
    \item igual a $1$, independientemente del valor de $\alpha + \beta$.
\end{opts}

\q{En una función de producción CES $Q = A[\delta K^\rho + (1-\delta)L^\rho]^{\nu/\rho}$, el parámetro $\nu$ determina:}
\begin{opts}
    \item la elasticidad de sustitución entre capital y trabajo ($\sigma$).
    \item el grado de homogeneidad (rendimientos a escala) de la tecnología.
    \item la participación constante del capital en el costo total de planta.
    \item el nivel mínimo de producción factible en el horizonte de corto plazo.
\end{opts}

\q{En una tecnología de proporciones fijas $Q = \min(aK, bL)$, la senda de expansión de largo plazo para cualesquiera precios positivos $(w, r)$:}
\begin{opts}
    \item coincide siempre con el rayo $aK = bL$, sin efecto sustitución.
    \item depende directamente de la relación de precios relativos $w/r$.
    \item se ubica sobre uno de los ejes mediante una solución de esquina.
    \item presenta una pendiente creciente conforme aumenta el nivel de $Q$.
\end{opts}

\q{Si una firma cuenta con tecnología de sustitutos perfectos $Q = aK + bL$ y enfrenta precios de factores $r$ y $w$, la minimización de costos produce una solución interior (usando ambos factores) únicamente cuando:}
\begin{opts}
    \item el precio del trabajo supera al producto marginal del capital.
    \item la elasticidad de sustitución toma un valor unitario ($\sigma = 1$).
    \item se cumple la igualdad exacta $b/a = w/r$.
    \item los rendimientos a escala son estrictamente decrecientes en planta.
\end{opts}

\q{En estudios empíricos de regulación de EPS de saneamiento (SUNASS), se suele preferir la función de costos o producción Translog frente a la Cobb-Douglas porque la Translog:}
\begin{opts}
    \item impone rendimientos constantes a escala en toda la red de agua.
    \item es una forma funcional flexible que permite elasticidades variables.
    \item garantiza que la elasticidad de sustitución sea siempre igual a $1$.
    \item requiere estimar menos parámetros econométricos que la Cobb-Douglas.
\end{opts}

\q{Una función de producción $Q = f(K,L)$ es homogénea de grado $k > 1$. Si la firma paga a cada factor exactamente su producto marginal ($w = p\cdot PMg_L$, $r = p\cdot PMg_K$), entonces por el teorema de Euler el pago total a los factores:}
\begin{opts}
    \item agota exactamente el valor total de la producción con beneficio nulo.
    \item es inferior al valor de la producción, dejando beneficios positivos.
    \item es independiente del grado de homogeneidad $k$ de la tecnología.
    \item supera el valor de la producción, generando pérdidas económicas.
\end{opts}

\q{¿Es compatible una función de producción que exhibe simultáneamente rendimientos crecientes a escala en el largo plazo y rendimientos marginales decrecientes en cada factor en el corto plazo?}
\begin{opts}
    \item Sí, como en $Q = K^{0.7}L^{0.7}$ donde $\alpha, \beta < 1$ y $\alpha+\beta > 1$.
    \item No, porque los rendimientos crecientes a escala exigen $PMg$ creciente siempre.
    \item Solo si la elasticidad de sustitución entre factores es exactamente igual a cero.
    \item No, porque la concavidad en cada factor implica rendimientos decrecientes a escala.
\end{opts}

\q{Para una tecnología diferenciable general (no necesariamente homogénea), la elasticidad de escala local $E(K,L)$ se expresa como:}
\begin{opts}
    \item el producto $\varepsilon_{Q,K} \cdot \varepsilon_{Q,L}$.
    \item la suma $\varepsilon_{Q,K} + \varepsilon_{Q,L}$.
    \item el cociente $\varepsilon_{Q,K} / \varepsilon_{Q,L}$.
    \item la diferencia $1 - (\varepsilon_{Q,K} + \varepsilon_{Q,L})$.
\end{opts}

\q{Si una función de producción es homotética, es decir $Q = F(g(K,L))$ con $F'>0$ y $g$ homogénea de grado $1$, entonces a precios de factores constantes su senda de expansión de largo plazo es:}
\begin{opts}
    \item una curva convexa cuya pendiente crece siempre con el nivel de $Q$.
    \item una recta únicamente si $F(\cdot)$ es lineal con rendimientos constantes.
    \item una línea recta desde el origen, aunque los rendimientos a escala varíen.
    \item horizontal o vertical dependiendo de qué insumo sea relativamente caro.
\end{opts}

\q{En una firma competitiva de largo plazo que posee una tecnología con rendimientos constantes a escala en todo el dominio y enfrenta precios dados $(p, w, r)$, el problema de maximización de beneficios tiene:}
\begin{opts}
    \item beneficio nulo si $p = CMeLP$, o solución no acotada si $p > CMeLP$.
    \item una única cantidad óptima positiva y finita para cualquier $p > CMeLP$.
    \item beneficios económicos estrictamente positivos en equilibrio competitivo.
    \item un costo medio de largo plazo creciente conforme aumenta la escala.
\end{opts}

\q{Si $f(K,L)$ es homogénea de grado $k$, entonces sus derivadas parciales $PMg_K$ y $PMg_L$ son funciones homogéneas de grado:}
\begin{opts}
    \item $k + 1$, creciendo más rápido que el producto total al escalar insumos.
    \item $k$, conservando exactamente el mismo grado que la función de producción.
    \item $0$ para cualquier valor del parámetro de homogeneidad $k > 0$.
    \item $k - 1$, por lo que bajo rendimientos constantes dependen solo de $K/L$.
\end{opts}

\q{En el problema de minimización de costos sujeto a producir $q$, la condición de tangencia interior $TMST_{L,K} = w/r$ equivale económicamente a que:}
\begin{opts}
    \item el gasto total en trabajo sea igual al gasto total en capital.
    \item el producto medio de ambos factores sea idéntico en el equilibrio.
    \item el último sol en trabajo rinde igual producto que en capital.
    \item el costo marginal de producción sea igual al salario nominal $w$.
\end{opts}

\q{En el Lagrangiano de minimización de costos $\mathcal{L} = wL + rK + \mu(q - f(K,L))$, el multiplicador $\mu^*$ evaluado en el óptimo representa:}
\begin{opts}
    \item el costo medio total de largo plazo $CMe(w,r,q) = C/q$.
    \item el costo marginal de producción $CMg(w,r,q)$.
    \item la tasa marginal de sustitución técnica en el equilibrio.
    \item el beneficio marginal obtenido por cada unidad adicional vendida.
\end{opts}

\q{La función de costos mínimos $C(w,r,q)$ es siempre cóncava respecto al vector de precios de los factores $(w,r)$ porque:}
\begin{opts}
    \item al cambiar precios, la firma sustituye y gasta menos que con insumos fijos.
    \item los rendimientos a escala de la firma deben ser decrecientes en largo plazo.
    \item el costo marginal de producción es siempre decreciente con el nivel de $q$.
    \item las isocuantas de producción deben ser estrictamente cóncavas al origen.
\end{opts}

\q{Respecto a los precios de los insumos $(w,r)$, la función de costos $C(w,r,q)$ es homogénea de grado $1$, mientras que las demandas condicionales de factores $L^c(w,r,q)$ y $K^c(w,r,q)$ son:}
\begin{opts}
    \item homogéneas de grado uno ($1$) en $(w,r)$.
    \item homogéneas de grado negativo ($-1$) en $(w,r)$.
    \item independientes de la razón de precios $w/r$.
    \item homogéneas de grado cero ($0$) en $(w,r)$.
\end{opts}

\q{¿Cuál es la relación matemática entre el costo total mínimo de largo plazo $C_{LP}(w,r,q)$ y el costo total de corto plazo $C_{CP}(w,r,q;\bar{K})$ para un mismo nivel de producción $q$?}
\begin{opts}
    \item $C_{LP}(w,r,q) \ge C_{CP}(w,r,q;\bar{K})$ porque el capital variable añade costos fijos.
    \item $C_{LP}(w,r,q) \le C_{CP}(w,r,q;\bar{K})$ para todo $q$, con igualdad cuando $\bar{K} = K^*(q)$.
    \item $C_{LP}(w,r,q) = C_{CP}(w,r,q;\bar{K})$ únicamente cuando existen rendimientos constantes.
    \item $C_{LP}(w,r,q) < C_{CP}(w,r,q;\bar{K})$ estrictamente para cualquier tamaño de planta $\bar{K}$.
\end{opts}

\q{Si una función de producción es homogénea de grado $k > 0$, entonces su función de costos de largo plazo $C(w,r,q)$ tiene la forma separable:}
\begin{opts}
    \item multiplicativa $C(w,r,q) = \phi(w,r) q^{k}$.
    \item aditiva $C(w,r,q) = \phi(w,r) + q^{1/k}$.
    \item producto $C(w,r,q) = \phi(w,r) q^{1/k}$.
    \item cociente $C(w,r,q) = \phi(w,r) / q^{k}$.
\end{opts}

\q{Las demandas condicionales de factores $L^c(w,r,q)$ y $K^c(w,r,q)$ cumplen necesariamente que $\frac{\partial L^c}{\partial w} \le 0$ porque:}
\begin{opts}
    \item no pueden existir insumos inferiores en la producción de corto plazo.
    \item el efecto producto siempre domina al efecto sustitución en factores.
    \item la función de producción debe ser estrictamente cóncava en trabajo.
    \item coincide con $\partial^2 C/\partial w^2 \le 0$ por concavidad de $C$.
\end{opts}

\q{La condición de segundo orden para que una combinación interior $(K^*, L^*)$ minimice costos sobre la isocuanta $f(K,L)=q$ exige que:}
\begin{opts}
    \item la isocuanta sea localmente convexa en el punto de tangencia.
    \item la función de producción sea estrictamente creciente a escala.
    \item el costo marginal sea estrictamente mayor que el costo medio.
    \item la función de producción sea necesariamente cóncava en $(K,L)$.
\end{opts}

\q{En teoría del productor, el lema de Shephard establece que si la función de costos $C(w,r,q)$ es diferenciable respecto al precio del insumo $w_i$, entonces $\frac{\partial C(w,r,q)}{\partial w_i}$ es igual a:}
\begin{opts}
    \item la demanda incondicional (maximizadora de beneficios) del factor $i$.
    \item la demanda condicional (minimizadora de costos) del factor $i$.
    \item el costo marginal de expandir el producto mediante el factor $i$.
    \item la participación del factor $i$ en el costo total de la empresa.
\end{opts}

\q{Por el teorema de Young aplicado a la función de costos $C(w,r,q)$, los efectos cruzados de precios sobre las demandas condicionales de factores satisfacen:}
\begin{opts}
    \item $\frac{\partial L^c}{\partial r} = -\frac{\partial K^c}{\partial w}$ (antisimetría de efectos).
    \item $\frac{\partial L^c}{\partial r} < 0$ siempre que haya solo dos factores.
    \item $\partial L^c/\partial r = \partial K^c/\partial w$ (simetría).
    \item $\frac{\partial L^c}{\partial r} = \frac{L^c}{K^c}$ en toda función homogénea.
\end{opts}

\q{La función de beneficios máximos $\pi(p, w, r)$ de una empresa precio-aceptante es siempre:}
\begin{opts}
    \item homogénea de grado $0$ y cóncava en el vector de precios $(p, w, r)$.
    \item homogénea de grado $1$ y cóncava en los precios de los insumos $(w,r)$.
    \item decreciente en el precio del producto $p$ y creciente en $(w,r)$.
    \item homogénea de grado $1$ y convexa en el vector de precios $(p, w, r)$.
\end{opts}

\q{Si la función de beneficios $\pi(p, w, r)$ es diferenciable, el lema de Hotelling permite recuperar directamente la función de oferta $q^*(p,w,r)$ y la demanda incondicional de trabajo $L^*(p,w,r)$ como:}
\begin{opts}
    \item $q^* = \frac{\partial \pi}{\partial p}$ y $L^* = -\frac{\partial \pi}{\partial w}$ (con signo menos).
    \item $q^* = -\frac{\partial \pi}{\partial p}$ (negativa) y $L^* = +\frac{\partial \pi}{\partial w}$ (positiva).
    \item ambas con signo positivo: $q^* = \frac{\partial \pi}{\partial p}$ y $L^* = \frac{\partial \pi}{\partial w}$.
    \item las segundas derivadas $q^* = \frac{\partial^2 \pi}{\partial p^2}$ y $L^* = -\frac{\partial^2 \pi}{\partial w^2}$.
\end{opts}

\q{De la convexidad de la función de beneficios $\pi(p,w,r)$ y el lema de Hotelling se deduce inmediatamente que para toda firma maximizadora de beneficios:}
\begin{opts}
    \item $\frac{\partial q^*}{\partial p} > 0$ solo si no existen insumos inferiores.
    \item $\partial q^*/\partial p \ge 0$ y $\partial L^*/\partial w \le 0$ siempre.
    \item $\frac{\partial L^*}{\partial w} > 0$ cuando el trabajo es un insumo inferior.
    \item $\frac{\partial q^*}{\partial p} \le 0$ en presencia de rendimientos decrecientes.
\end{opts}

\q{Para que exista un máximo interior finito de beneficios en una firma precio-aceptante, la tecnología debe exhibir localmente en el óptimo:}
\begin{opts}
    \item rendimientos crecientes a escala con costo marginal decreciente.
    \item rendimientos constantes a escala con $p$ estrictamente mayor que $CMg$.
    \item rendimientos decrecientes a escala (función $f$ localmente cóncava).
    \item elasticidad de sustitución técnica igual a cero entre insumos.
\end{opts}

\q{En una firma precio-aceptante que maximiza beneficios, la condición de primer orden para contratar trabajo $L$ establece que $p \cdot PMg_L = w$, lo cual significa que:}
\begin{opts}
    \item el valor del producto marginal del trabajo iguala al salario $w$.
    \item el producto medio del trabajo iguala al salario real $w/p$.
    \item el costo medio de producción iguala al ingreso marginal del bien.
    \item la elasticidad del producto respecto al trabajo es igual a $1$.
\end{opts}

\q{A diferencia de la teoría del consumidor donde la demanda marshalliana puede tener pendiente positiva (bien Giffen), la demanda incondicional de un factor $L^*(p,w,r)$ nunca puede tener pendiente positiva respecto a $w$ porque:}
\begin{opts}
    \item todos los factores productivos son necesariamente sustitutos netos.
    \item el producto marginal de los insumos es siempre constante en $L$.
    \item los insumos inferiores no pueden existir bajo minimización de costos.
    \item la empresa no enfrenta una restricción presupuestaria fija.
\end{opts}

\q{En las curvas de costos de corto plazo en forma de U, la curva de costo marginal ($CMg$) intercepta a las curvas de costo variable medio ($CVMe$) y costo medio total ($CMe$):}
\begin{opts}
    \item exactamente en sus puntos máximos, de arriba hacia abajo.
    \item en el punto donde el costo fijo medio ($CFMe$) iguala al $CVMe$.
    \item en sus respectivos mínimos, de abajo hacia arriba.
    \item únicamente en el nivel de producción donde $PMg_L$ es máximo.
\end{opts}

\q{El nivel de producción $q_1$ donde $CVMe$ alcanza su mínimo es siempre menor que el nivel $q_2$ donde $CMe$ alcanza su mínimo cuando existen costos fijos positivos ($CF > 0$) porque:}
\begin{opts}
    \item el $CMg$ es estrictamente decreciente entre $q_1$ y $q_2$.
    \item el $CFMe$ sigue cayendo en $q_1$, reduciendo aún al $CMe$.
    \item el $CF$ crece proporcionalmente al producto entre $q_1$ y $q_2$.
    \item en $q_1$ se cumple que $CMe = CVMe$ y luego $CMe$ desciende.
\end{opts}

\q{Cuando el trabajo $L$ es el único insumo variable en el corto plazo con salario constante $w$, la relación exacta entre el costo marginal de corto plazo ($CMg$) y el producto marginal ($PMg_L$) es:}
\begin{opts}
    \item $CMg = w/PMg_L$, siendo mínimo cuando $PMg_L$ es máximo.
    \item $CMg = w \cdot PMg_L$, moviéndose ambos en la misma dirección.
    \item $CMg = \frac{w}{PMe_L}$, alcanzando su mínimo en el máximo de $PMe_L$.
    \item $CMg = \frac{PMg_L}{w}$, siendo independiente del nivel de producto.
\end{opts}

\q{En el corto plazo, si los costos fijos $CF$ son completamente hundidos (irrecuperables), una empresa competitiva decide cerrar ($q^* = 0$) cuando el precio de mercado $p$ satisface:}
\begin{opts}
    \item $p < \min CMe$, pues toda pérdida económica obliga al cierre inmediato.
    \item $p < CFMe$, porque los ingresos no cubren los costos hundidos.
    \item $\min CVMe < p < \min CMe$, al no recuperarse la inversión fija.
    \item $p < \min CVMe$, pues producir generaría pérdidas mayores que $CF$.
\end{opts}

\q{Si una empresa enfrenta un costo fijo irrecuperable $CF_h$ y además un costo cuasi-fijo $F_c > 0$ que solo se paga si $q > 0$ (evitable al cerrar), su regla óptima de cierre a corto plazo es cerrar cuando:}
\begin{opts}
    \item $p < \min CVMe(q)$, ignorando por completo el costo cuasi-fijo $F_c$.
    \item $p < \min \left(\frac{CV(q) + F_c}{q}\right)$ (costo medio evitable).
    \item $p < \min \left(\frac{CV(q) + CF_h}{q}\right)$ (ignorando $F_c$).
    \item $p < \min CMe(q)$, incluyendo tanto $CF_h$ como $F_c$ en el umbral.
\end{opts}

\q{Para una firma competitiva con costos fijos totalmente hundidos, su curva de oferta individual de corto plazo coincide con:}
\begin{opts}
    \item toda la curva de $CMg$ para cualquier nivel de precio positivo.
    \item el tramo creciente del $CMg$ situado por encima del mínimo de $CMe$.
    \item el tramo creciente del $CMg$ por encima del mínimo de $CVMe$.
    \item el tramo decreciente del $CMe$ hasta el punto de eficiencia mínima.
\end{opts}

\q{En economía de la regulación de servicios públicos (como infraestructura enterrada de redes de agua potable y alcantarillado), un costo hundido se caracteriza por:}
\begin{opts}
    \item variar proporcionalmente con los metros cúbicos de agua facturados.
    \item ser totalmente recuperable mediante la reventa inmediata de los activos.
    \item determinar el nivel óptimo de producción en el margen a corto plazo.
    \item carecer de uso alternativo y no poder recuperarse al salir.
\end{opts}

\q{La curva de costo medio de largo plazo ($CMeLP$) es la envolvente inferior de las curvas de costo medio de corto plazo ($CMeCP$). En un punto de tangencia situado en el tramo decreciente de $CMeLP$, la planta elegida opera:}
\begin{opts}
    \item a la izquierda de su mínimo de $CMeCP$ (con capacidad ociosa).
    \item exactamente en el punto mínimo de su curva $CMeCP$ de corto plazo.
    \item a la derecha de su mínimo de $CMeCP$ (sobreutilizando la planta).
    \item en un punto donde $CMgLP$ supera estrictamente a $CMeLP$.
\end{opts}

\q{En el nivel de producción $q_0$ donde una curva de costo medio de corto plazo $CMeCP(q; \bar{K})$ es tangente a la envolvente $CMeLP(q)$, se cumple además que:}
\begin{opts}
    \item desigualdad estricta: $CMgCP(q_0; \bar{K}) > CMgLP(q_0)$.
    \item desigualdad inversa: $CMgCP(q_0; \bar{K}) < CMgLP(q_0)$.
    \item igualdad exacta: $CMgCP(q_0; \bar{K}) = CMgLP(q_0)$.
    \item coincidencia con el medio: $CMgCP(q_0; \bar{K}) = CMeLP(q_0)$.
\end{opts}

\q{La elasticidad del costo total de largo plazo respecto al producto, $\varepsilon_{C,q} = \frac{CMgLP}{CMeLP}$, permite identificar las economías de escala de la empresa. Existen economías de escala cuando:}
\begin{opts}
    \item $\varepsilon_{C,q} > 1$, lo cual implica $CMgLP > CMeLP$ y $CMeLP$ creciente.
    \item $\varepsilon_{C,q} < 1$, implicando $CMgLP < CMeLP$ y $CMeLP$ cayendo.
    \item $\varepsilon_{C,q} = 1$, lo cual implica $CMgLP = CMeLP$ en todo el dominio.
    \item $\varepsilon_{C,q} < 0$, lo cual implica costos totales decrecientes con $q$.
\end{opts}

\q{La escala mínima eficiente (EME) de una industria se define como:}
\begin{opts}
    \item el nivel máximo de producción que una planta puede generar en corto plazo.
    \item la cantidad producida donde el costo marginal corta al costo variable medio.
    \item el volumen mínimo requerido para que el costo fijo medio sea cero.
    \item el menor nivel de producción donde $CMeLP$ alcanza su mínimo.
\end{opts}

\q{Comparando las pendientes de las curvas de costo marginal en el punto de tangencia $q_0$ donde $CMgCP(q_0) = CMgLP(q_0)$, la curva de $CMgCP$ es siempre:}
\begin{opts}
    \item más empinada que $CMgLP$ por la rigidez del capital.
    \item más plana que $CMgLP$ porque en corto plazo no se paga capital nuevo.
    \item idéntica en pendiente a $CMgLP$ por el teorema de la envolvente.
    \item horizontal en $q_0$ independientemente de los rendimientos a escala.
\end{opts}

\q{El principio de Le Chatelier aplicado a la teoría del productor implica que la curva de oferta de largo plazo de una empresa competitiva es:}
\begin{opts}
    \item menos elástica que su curva de oferta de corto plazo ante cambios de $p$.
    \item más elástica que la oferta de corto plazo en el equilibrio.
    \item perfectamente inelástica cuando todos los factores son ajustables.
    \item de pendiente negativa siempre que existan insumos inferiores.
\end{opts}

\q{Un factor productivo $L$ se clasifica como un insumo inferior (o regresivo) en un tramo de producción cuando, a precios de factores constantes:}
\begin{opts}
    \item $\frac{\partial L^c(w,r,q)}{\partial w} > 0$ (su demanda sube al subir $w$).
    \item $PMg_L < 0$ en el punto donde se minimizan los costos de planta.
    \item $\partial L^c/\partial q < 0$ (su uso cae al expandirse $q$).
    \item su participación en el costo total decrece al aumentar el salario $w$.
\end{opts}

\q{Si el trabajo $L$ es un insumo inferior ($\frac{\partial L^c}{\partial q} < 0$), entonces un incremento en el salario $w$ provoca que el costo marginal de producción $CMg(w,r,q)$:}
\begin{opts}
    \item disminuya, desplazando hacia abajo la curva de $CMg$ y elevando $q^*$.
    \item aumente necesariamente en una proporción superior al aumento de $w$.
    \item permanezca inalterado porque el efecto sustitución anula al de escala.
    \item se vuelva negativo en el tramo donde el trabajo es regresivo.
\end{opts}

\q{Al descomponer el efecto total de un alza en el salario $w$ sobre la demanda incondicional de trabajo $L^*(p,w,r)$ en efecto sustitución y efecto producto (o escala), para un insumo normal ambos efectos:}
\begin{opts}
    \item tienen signos opuestos, pudiendo el efecto producto dominar al sustitución.
    \item son positivos porque la firma expande producción para cubrir costos.
    \item se cancelan mutuamente si la tecnología tiene rendimientos constantes.
    \item son negativos y se refuerzan, cayendo $L^*$ más que $L^c$.
\end{opts}

\q{En una empresa que emplea únicamente dos insumos ($K$ y $L$) con isocuantas suaves y estrictamente convexas, los dos factores deben ser necesariamente:}
\begin{opts}
    \item complementos técnicos netos ($\frac{\partial L^c}{\partial r} < 0$).
    \item sustitutos técnicos netos ($\frac{\partial L^c}{\partial r} > 0$).
    \item complementos brutos ($\frac{\partial L^*}{\partial r} < 0$) en todo caso.
    \item insumos independientes con efecto cruzado nulo ($\frac{\partial L^c}{\partial r} = 0$).
\end{opts}

\q{Una función de costos $C(q)$ es estrictamente subaditiva en el nivel de producto $q$ (condición definitoria de monopolio natural) cuando:}
\begin{opts}
    \item el costo marginal de largo plazo es estrictamente menor que cero en $q$.
    \item el costo medio de largo plazo es decreciente en todo el intervalo $(0, \infty)$.
    \item una firma produce $q$ más barato que cualquier partición en varias.
    \item la firma obtiene beneficios económicos positivos al cobrar el costo marginal.
\end{opts}

\q{En el sector saneamiento regulado por SUNASS, una EPS presta conjuntamente los servicios de agua potable ($q_1$) y alcantarillado ($q_2$). Existen economías de alcance (o de ámbito) cuando:}
\begin{opts}
    \item el costo conjunto es mayor: $C(q_1, q_2) > C(q_1, 0) + C(0, q_2)$.
    \item existe homogeneidad de escala: $C(tq_1, tq_2) < t C(q_1, q_2)$ para $t > 1$.
    \item la suma de costos marginales cumple $\frac{\partial C}{\partial q_1} + \frac{\partial C}{\partial q_2} < C(q_1, q_2)$.
    \item el costo conjunto es menor: $C(q_1, q_2) < C(q_1, 0) + C(0, q_2)$.
\end{opts}

\q{En una firma multiproducto como una EPS (agua potable y alcantarillado), ¿son las economías de escala específicas por producto suficientes para garantizar la subaditividad de costos?}
\begin{opts}
    \item No; también se requieren economías de alcance o complementariedad de costos.
    \item Sí, las economías de escala individuales garantizan siempre la subaditividad.
    \item No, porque la subaditividad multiproducto solo depende de los costos hundidos.
    \item Sí, siempre que los costos fijos sean superiores a los costos variables.
\end{opts}

\q{En la industria de agua potable y saneamiento (SUNASS), las economías de densidad de red se manifiestan cuando:}
\begin{opts}
    \item el costo total aumenta proporcionalmente al número de distritos atendidos.
    \item la EPS expande su red hacia zonas rurales dispersas reduciendo el $CMe$.
    \item el costo medio por usuario cae al atender más conexiones en una misma red.
    \item la tarifa media del primer bloque iguala al costo marginal de bombeo.
\end{opts}

\q{En el caso uniproducto, ¿cuál es la relación lógica correcta entre economías de escala globales ($CMeLP$ decreciente en todo el rango) y subaditividad de costos?}
\begin{opts}
    \item Las economías de escala globales son condición necesaria y suficiente para la subaditividad.
    \item Son condición suficiente, pero no necesaria, para la subaditividad.
    \item La subaditividad es condición suficiente, pero no necesaria, para economías de escala globales.
    \item No existe ninguna relación matemática entre $CMeLP$ decreciente y subaditividad.
\end{opts}

\q{Si una EPS que opera como monopolio natural con $CMg < CMe$ fijara su tarifa exactamente igual al costo marginal (eficiencia asignativa de primer mejor), entonces:}
\begin{opts}
    \item sufriría déficit, requiriendo subsidio o tarifa en dos partes.
    \item obtendría beneficios extraordinarios que SUNASS debería expropiar.
    \item cubriría exactamente sus costos totales medios con beneficio nulo.
    \item perdería las economías de alcance entre agua potable y alcantarillado.
\end{opts}

\q{El progreso tecnológico se denomina neutral en el sentido de Hicks cuando, para una relación capital-trabajo ($K/L$) dada:}
\begin{opts}
    \item la razón capital-producto ($K/Q$) permanece constante a una tasa $r$ fija.
    \item la razón trabajo-producto ($L/Q$) permanece constante a un salario $w$ fijo.
    \item el producto marginal del capital aumenta mientras el del trabajo cae.
    \item la $TMST_{L,K}$ permanece constante al desplazarse las isocuantas.
\end{opts}

\q{El progreso tecnológico aumentador de trabajo, expresado como $Q = F(K, A(t)L)$, corresponde a la definición de progreso tecnológico neutral en el sentido de:}
\begin{opts}
    \item Solow (aumentador puro de capital $Q = F(A(t)K, L)$).
    \item Hicks (desplaza la función multiplicativamente por $A(t)$).
    \item Harrod (mantiene $K/Q$ constante si $PMg_K$ no varía).
    \item Leontief (reduce la elasticidad de sustitución hacia cero).
\end{opts}

\q{En una función de producción Cobb-Douglas $Q = A K^\alpha L^{1-\alpha}$, las definiciones de progreso tecnológico neutral de Hicks, Harrod y Solow:}
\begin{opts}
    \item son mutuamente excluyentes salvo cuando $\alpha = 0.5$.
    \item son matemáticamente equivalentes entre sí.
    \item difieren en el signo de su impacto sobre la $TMST_{L,K}$.
    \item solo coinciden si la elasticidad de sustitución es infinita.
\end{opts}

\q{En la contabilidad del crecimiento (residuo de Solow) bajo rendimientos constantes a escala y mercados competitivos, la tasa de crecimiento de la Productividad Total de Factores ($\hat{A}$) se calcula como:}
\begin{opts}
    \item la suma ponderada $\hat{A} = \hat{Q} + s_K \hat{K} + s_L \hat{L}$.
    \item la diferencia simple $\hat{A} = \hat{Q} - \hat{K} - \hat{L}$.
    \item el aporte factorial $\hat{A} = s_K \hat{K} + s_L \hat{L}$.
    \item el residuo $\hat{A} = \hat{Q} - s_K \hat{K} - s_L \hat{L}$.
\end{opts}

\subsection*{Parte II. Aplicaciones prácticas y ejercicios cuantitativos (154 a 180)}

\q{Una empresa tiene la función de producción de corto plazo $Q(L) = 18L^2 - L^3$. ¿A partir de qué nivel de trabajo $L$ comienza la etapa II de producción (donde $PMe_L$ alcanza su máximo)?}
\begin{opts}
    \item $L = 9.0$
    \item $L = 6.0$
    \item $L = 12.0$
    \item $L = 18.0$
\end{opts}

\q{Para la función de producción de corto plazo $Q(L) = 18L^2 - L^3$, el producto total alcanza su máximo (inicio de la etapa III) en $L = 12$, siendo el producto máximo igual a:}
\begin{opts}
    \item $729$ unidades
    \item $864$ unidades
    \item $972$ unidades
    \item $648$ unidades
\end{opts}

\q{Dada la función de producción $Q(L) = 20L^2 - L^3$, calcule la elasticidad del producto respecto al trabajo ($\varepsilon_{Q,L}$) evaluada en $L = 12$:}
\begin{opts}
    \item $0.75$
    \item $1.00$
    \item $0.50$
    \item $1.25$
\end{opts}

\q{Para la tecnología $Q = 5K^{0.25}L^{0.75}$, la tasa marginal de sustitución técnica del trabajo por capital ($TMST_{L,K}$) evaluada en $(K, L) = (16, 8)$ es:}
\begin{opts}
    \item $6.0$
    \item $3.0$
    \item $4.0$
    \item $8.0$
\end{opts}

\q{Sea la función de producción CES $Q = \left(0.4K^{-0.25} + 0.6L^{-0.25}\right)^{-4}$. ¿Cuál es el valor de su elasticidad de sustitución técnica $\sigma$?}
\begin{opts}
    \item $0.25$
    \item $1.25$
    \item $4.00$
    \item $0.80$
\end{opts}

\q{Determine el grado de homogeneidad y el tipo de rendimientos a escala de la tecnología $Q(K,L) = 2(K^{0.5} + L^{0.5})^3$:}
\begin{opts}
    \item grado $3.0$ (rendimientos crecientes a escala).
    \item grado $1.0$ (rendimientos constantes a escala).
    \item grado $1.5$ (rendimientos crecientes a escala).
    \item grado $0.5$ (rendimientos decrecientes a escala).
\end{opts}

\q{Para la función de producción no homogénea $Q(K,L) = KL + 2K$, calcule la elasticidad de escala local $E(K,L)$ evaluada en $(K, L) = (2, 3)$:}
\begin{opts}
    \item $1.20$
    \item $1.60$
    \item $1.80$
    \item $2.00$
\end{opts}

\q{Una empresa produce con $Q = K^{0.5}L^{0.5}$ y enfrenta precios de factores $w = 4$ y $r = 16$. El costo mínimo para producir $q = 20$ unidades es:}
\begin{opts}
    \item $S/\ 320.00$
    \item $S/\ 240.00$
    \item $S/\ 400.00$
    \item $S/\ 160.00$
\end{opts}

\q{Si una planta de tratamiento tiene tecnología de proporciones fijas $Q = \min(2K, 5L)$ con precios de insumos $w = 10$ y $r = 6$, ¿cuál es el costo mínimo para producir $q = 40$?}
\begin{opts}
    \item $S/\ 160.00$
    \item $S/\ 240.00$
    \item $S/\ 320.00$
    \item $S/\ 200.00$
\end{opts}

\q{Para la función de producción $Q = 3K + 2L$ con $w = 6$ y $r = 12$, la combinación óptima $(K^*, L^*)$ y el costo mínimo para producir $q = 60$ son:}
\begin{opts}
    \item $(20, 0)$ con costo $S/\ 240$
    \item $(0, 30)$ con costo $S/\ 180$
    \item $(10, 15)$ con costo $S/\ 210$
    \item $(0, 20)$ con costo $S/\ 120$
\end{opts}

\q{Dada la tecnología CES $Q = (K^{0.5} + L^{0.5})^2$ con precios de factores $w = 1$ y $r = 4$, el costo total mínimo para producir $q = 25$ es:}
\begin{opts}
    \item $S/\ 25.00$
    \item $S/\ 30.00$
    \item $S/\ 20.00$
    \item $S/\ 15.00$
\end{opts}

\q{Una empresa con tecnología $Q = K^{0.5}L$ enfrenta $w = 8$ y $r = 4$. Para producir $q = 8$ unidades al mínimo costo, la cantidad óptima de trabajo $L^c$ y el costo mínimo son:}
\begin{opts}
    \item $L^c = 4$ y $C = S/\ 48$
    \item $L^c = 2$ y $C = S/\ 32$
    \item $L^c = 8$ y $C = S/\ 64$
    \item $L^c = 6$ y $C = S/\ 56$
\end{opts}

\q{La función de costos de largo plazo de una empresa es $C(w, r, q) = 4q w^{0.25}r^{0.75}$. Aplicando el lema de Shephard, la demanda condicional de trabajo $L^c$ cuando $w = 1$, $r = 16$ y $q = 10$ es:}
\begin{opts}
    \item $40.0$ unidades
    \item $60.0$ unidades
    \item $20.0$ unidades
    \item $80.0$ unidades
\end{opts}

\q{Sea la función de costos $C(w, r, q) = q^2(2w + 3r + 2\sqrt{wr})$. Evaluada en $w = 9$, $r = 4$ y $q = 2$, la demanda condicional de capital $K^c$ es:}
\begin{opts}
    \item $12.0$ unidades
    \item $18.0$ unidades
    \item $15.0$ unidades
    \item $24.0$ unidades
\end{opts}

\q{Si la función de costos de largo plazo es $C(w, r, q) = 3q^2 w^{0.5}r^{0.5}$, ¿cuál es el costo marginal $CMg$ cuando $w = 4$, $r = 9$ y $q = 5$?}
\begin{opts}
    \item $S/\ 150.00$
    \item $S/\ 210.00$
    \item $S/\ 180.00$
    \item $S/\ 90.00$
\end{opts}

\q{Una empresa competitiva produce en el corto plazo con $Q = 10\sqrt{L}$, enfrenta $p = 6$, $w = 15$ y un costo fijo $CF = 40$. El trabajo óptimo $L^*$ y el beneficio máximo $\pi^*$ son:}
\begin{opts}
    \item $L^* = 4$ y $\pi^* = S/\ 20$
    \item $L^* = 4$ y $\pi^* = S/\ 60$
    \item $L^* = 9$ y $\pi^* = S/\ 25$
    \item $L^* = 2$ y $\pi^* = S/\ 15$
\end{opts}

\q{Sea una empresa precio-aceptante con tecnología $Q = 4K^{0.25}L^{0.25}$ que enfrenta $p = 8$, $w = 2$ y $r = 2$. El nivel de producción óptimo $Q^*$ y el beneficio máximo $\pi^*$ son:}
\begin{opts}
    \item $Q^* = 16$ y $\pi^* = S/\ 96$
    \item $Q^* = 16$ y $\pi^* = S/\ 64$
    \item $Q^* = 32$ y $\pi^* = S/\ 64$
    \item $Q^* = 8$ y $\pi^* = S/\ 32$
\end{opts}

\q{La función de beneficios de una firma competitiva es $\pi(p, w, r) = \frac{5p^2}{\sqrt{wr}}$. Usando el lema de Hotelling, la cantidad ofrecida $q^*$ cuando $p = 18$, $w = 4$ y $r = 9$ es:}
\begin{opts}
    \item $15.0$ unidades
    \item $45.0$ unidades
    \item $60.0$ unidades
    \item $30.0$ unidades
\end{opts}

\q{Dada la función de beneficios $\pi(p, w, r) = \frac{5p^2}{w^{0.5}r^{0.5}}$, halle mediante el lema de Hotelling la demanda incondicional de trabajo $L^*$ en $p = 12$, $w = 4$ y $r = 9$:}
\begin{opts}
    \item $30.0$ unidades
    \item $20.0$ unidades
    \item $15.0$ unidades
    \item $10.0$ unidades
\end{opts}

\q{Una empresa competitiva tiene la función de costos totales de corto plazo $CT(q) = 50 + 20q - 4q^2 + 0.5q^3$ (con costo fijo hundido). El precio de cierre por debajo del cual la firma produce $q^* = 0$ es:}
\begin{opts}
    \item $S/\ 16.00$
    \item $S/\ 12.00$
    \item $S/\ 20.00$
    \item $S/\ 10.00$
\end{opts}

\q{Para la función de costos $CT(q) = 100 + 10q + q^2$, el nivel de producción que minimiza el costo medio total ($CMe$) y dicho costo medio mínimo son:}
\begin{opts}
    \item $q = 10$ y $CMe_{\min} = S/\ 30$
    \item $q = 10$ y $CMe_{\min} = S/\ 20$
    \item $q = 5$ y $CMe_{\min} = S/\ 35$
    \item $q = 20$ y $CMe_{\min} = S/\ 35$
\end{opts}

\q{Una empresa tiene costos hundidos de $S/\ 60$, costos cuasi-fijos evitables de $S/\ 16$ (si $q > 0$) y costos variables $CV(q) = 4q + q^2$. ¿Cuál es el precio mínimo a partir del cual decide operar ($q > 0$)?}
\begin{opts}
    \item $S/\ 4.00$
    \item $S/\ 16.00$
    \item $S/\ 12.00$
    \item $S/\ 20.00$
\end{opts}

\q{Una empresa enfrenta costos de corto plazo $C_{CP}(q, K) = 4K + \frac{q^2}{K}$. En el largo plazo (eligiendo óptimamente el tamaño de planta $K$), el costo total mínimo para producir $q = 30$ es:}
\begin{opts}
    \item $S/\ 150.00$
    \item $S/\ 180.00$
    \item $S/\ 90.00$
    \item $S/\ 120.00$
\end{opts}

\q{La función de costos totales de largo plazo de una industria es $C_{LP}(q) = q^3 - 12q^2 + 60q$. La escala mínima eficiente (EME) y el costo medio mínimo correspondiente son:}
\begin{opts}
    \item $q = 4$ y $CMeLP = S/\ 28$
    \item $q = 6$ y $CMeLP = S/\ 24$
    \item $q = 8$ y $CMeLP = S/\ 28$
    \item $q = 6$ y $CMeLP = S/\ 36$
\end{opts}

\q{Una EPS de agua potable presenta costos totales $C(q) = 100 + 10q + 0.5q^2$ (en miles de soles). Evaluada en $q = 10$, la elasticidad costo-producto ($\varepsilon_{C,q}$) es:}
\begin{opts}
    \item $1.25$ (deseconomías de escala)
    \item $1.00$ (rendimientos constantes)
    \item $0.80$ (economías de escala)
    \item $0.50$ (economías de escala)
\end{opts}

\q{Una EPS brinda agua potable ($q_1$) y alcantarillado ($q_2$) con una función de costos $C(q_1, q_2) = 100 + 5q_1 + 5q_2$ (donde $100$ es un costo fijo común que se paga si se produce al menos un servicio). Para $q_1 = 30$ y $q_2 = 30$, el grado de economías de alcance $S_C$ es:}
\begin{opts}
    \item $0.20$
    \item $0.50$
    \item $0.40$
    \item $0.25$
\end{opts}

\q{Una empresa de saneamiento tiene costos totales $C(q) = 36 + q^2$ (cuyo $CMe$ mínimo ocurre en $q = 6$). Si la demanda total del mercado es $q = 8$, compare el costo de un solo operador frente a dos operadores simétricos ($q_1 = q_2 = 4$):}
\begin{opts}
    \item es subaditiva ($S/\ 100$ con uno frente a $S/\ 104$ con dos).
    \item no es subaditiva ($S/\ 104$ con uno frente a $S/\ 100$ con dos).
    \item es indiferente ($S/\ 100$ tanto con uno como con dos operadores).
    \item no es subaditiva ($S/\ 100$ con uno frente a $S/\ 72$ con dos).
\end{opts}
